% !TeX root = ../../Book2.tex
\section{与限制问题的关系}
\label{b2:sec:2404}

球面限制控制薄壳上的 Fourier 质量，Bochner–Riesz 符号则把不同薄壳加权求和。两者共享曲率与频率集中，但混合指数界和同指数乘子界不同，不能仅凭一个球面定理直接宣称整个乘子猜想成立。

\subsection{球面 Fourier 变换}
\label{b2:subsec:240401}

极坐标给出层分解
\[
 B_1^\delta f(x)=\int_0^1(1-r^2)^\delta r^{d-1}
        \int_{S^{d-1}}e^{2\pi\mathrm i r x\cdot\omega}
                  \widehat f(r\omega)\,\mathrm d\sigma(\omega)\,\mathrm dr.
\]
在 Schwartz 输入上换序合法，因为球内符号可积。内层正是半径 \(r\) 的球面限制再延拓。
球面本身的 Fourier 变换在径向尾部衰减，径向权的边界消失再增加 \(\delta+1\) 的振荡增益；这与核渐近的总阶数一致。

\ProseParagraphBreak
若只把内层绝对值估计后对 \(r\) 积分，就会丢失不同半径之间的抵消。某些混合指数估计仍可由这一办法直接证明，但锐利同指数乘子问题往往需要保留空间与径向的振荡。

\subsection{限制估计}
\label{b2:subsec:240402}

令 \(p_*=2(d+1)/(d+3)\)、\(q_*=p_*'\)。对薄壳
\(\Omega_\epsilon=\{1-\epsilon<|\xi|<1+\epsilon\}\)，球面限制的缩放给出
\[
 \|\mathbf1_{\Omega_\epsilon}\widehat f\|_2
       \le C\epsilon^{1/2}\|f\|_{p_*},\qquad0<\epsilon<1/2.
\]
证明是将平方范数写成半径积分，逐半径应用限制界；半径在固定紧区间，其缩放常数一致。
与对偶估计复合，或直接逐半径限制再延拓，得到
\[
 \|\mathcal F^{-1}(\mathbf1_{\Omega_\epsilon}\widehat f)\|_{q_*}
       \le C\epsilon\|f\|_{p_*}.
\]
第一式控制一个频率 \(L^2\) 范数；第二式是 \(L^{p_*}\to L^{q_*}\) 的混合指数界，不是 \(L^{p_*}\to L^{p_*}\)。

\ProseParagraphBreak
对薄壳内一个非负径向权 \(w(r)\)，同样有算子范数至多
\(C\int|w(r)|\,\mathrm dr\)。在边界距离约 \(2^{-j}\) 的壳层上，Bochner–Riesz 权大小约 \(2^{-j\delta}\)，得到混合指数贡献 \(C2^{-j(\delta+1)}\)。因此在这条混合指数线上，壳层求和对 \(\delta>-1\) 可收敛；这与负参数不属于 \(L^2\) 有界乘子并不矛盾。

把半径趋零的缩放也保留下来，可以完整证明一个全平均的混合指数界。
\begin{proposition}\label{b2:prop:24-mixed-bound}
设 \(d\geq2\)、\(p_*=2(d+1)/(d+3)\)、\(q_*=p_*'\)。对任意 \(\delta>-1\)，以球内可积符号在 Schwartz 类上定义的 \(B_1^\delta\) 有唯一连续延拓
\[
 B_1^\delta:L^{p_*}(\mathbb R^d)\longrightarrow L^{q_*}(\mathbb R^d).
\]
\end{proposition}
\begin{proof}
记 \(T_r f(x)=\int_{S^{d-1}}e^{2\pi\mathrm irx\cdot\omega}
\widehat f(r\omega)\,\mathrm d\sigma(\omega)\)。设 \(f_r(y)=f(y/r)\)，则
\(\widehat f(r\omega)=r^{-d}\widehat{f_r}(\omega)\)，因而
\[
 T_r f(x)=r^{-d}(ERf_r)(rx),\qquad
 \|T_r f\|_{q_*}\leq C r^{d/p_*-d/q_*-d}\|f\|_{p_*}.
\]
层分解中的体积权 \(r^{d-1}\) 使最后指数成为
\(d/p_*-d/q_*-1=(d-1)/(d+1)>-1\)。
于是 Minkowski 与限制端点给出
\[
 \|B_1^\delta f\|_{q_*}\leq C\|f\|_{p_*}
 \int_0^1(1-r^2)^\delta r^{(d-1)/(d+1)}\,\mathrm dr.
\]
积分在 \(r=0\) 的幂次可积，在 \(r=1\) 可积恰需 \(\delta>-1\)。稠密性完成延拓。
\end{proof}
这条命题包括负参数，但改变了输出指数，因此不与负参数符号在边界无界矛盾。对任意有界同指数 \(L^2\) 乘子，Plancherel 需要符号本性有界；这里使用的是较强可积性输入换取较弱的全空间输出，根本没有提出平方到平方的界。

\ProseParagraphBreak

半径 \(R\) 的相同混合界还会带缩放因子
\(R^{d(1/p_*-1/q_*)}\)，不像同指数乘子范数那样独立于 \(R\)。如果在转述中删除这个因子，便会把不同类型的尺度关系混在一起。对固定单位半径的估计，缩放检查仍然有价值：它可以迅速发现体积权、曲面测度与 Fourier 归一化中的遗漏。

\subsection{乘子估计}
\label{b2:subsec:240403}

若要求同指数 \(L^p\)，上述输出 \(q_*>p_*\) 不能通过无限测度空间上的包含关系改成 \(p_*\)。一种方法是在空间中再按尺度局部化核，使用限制界控制一部分项、核尾部控制另一部分项，然后在壳层间建立可求和的界。这个过程涉及局部球的体积因子与几乎正交，不是形式上的指数替换。

\ProseParagraphBreak
\chapref{b2:ch:19}平方函数处理互相分离的频率环带，而本章的壳层都贴近同一个球面。它们在频率上可不交，却仍可能在空间中沿相近方向高度相干。普通 Plancherel 正交只给 \(p=2\)，其他指数上的组合需额外曲率或多线性控制。
现代解耦和限制方法正是组织这种相干的重要工具，具体结果必须按已出版定理核对。

\subsection{曲率}
\label{b2:subsec:240404}

球、椭球与平面多面体提供三种可比较几何。固定非退化线性换元把椭球变成球，因此同指数有界性范围及算子范数相同；各方向导数的估计则可能依赖换元矩阵。矩形可拆成各坐标区间投影，硬截断在 \(1<p<\infty\) 上有界。球面法向连续变化，硬截断则受 Fefferman 定理阻碍。

\ProseParagraphBreak
曲率既能改善单个振荡核衰减，也能产生大量不同方向的集中管。前一种效应使正平滑平均有较好的核估计，后一种效应使硬球乘子失败。它们并不矛盾，因为核尾部的单输入检验和多方向向量值检验衡量不同机制。
阅读乘子范围时，应明确平滑量、输入是否径向、空间是否局部、半径是否统一及是否有尺度损失；这些条件决定一个结论能否被用于普通 \(L^p(\mathbb R^d)\) 平均。

曲率条件的比较还可以作成实际的变换。若 \(A\) 为可逆矩阵，椭球符号 \((1-|A\xi|^2)_+^\delta\) 的算子由 Fourier 换元与空间线性拉回共轭于球平均。拉回在 \(L^p\) 的范数因子为 \(|\det A|^{1/p}\) 或其逆，共轭时两者相消；同指数算子范数因此与标准球相同。若改变的同时还使用不同的空间欧氏导数范数，其正则性常数才会依赖矩阵的条件数。乘子有界性与导数估计的矩阵依赖不能混为一谈。

\ProseParagraphBreak

多面体的情况并非只由“没有曲率”决定。一个固定有界凸多面体是有限个半空间的交，每个半空间乘子可通过旋转、调制和一维 Hilbert 投影在 \(1<p<\infty\) 上实现；这些 Fourier 乘法算子彼此交换，有限次复合就得到多面体指示函数乘子。其界可能随面数积累，所以用越来越多面的多面体逼近球，并没有产生球乘子的一致界。几何对象的逐点逼近若缺少一致算子常数，不能用于转移有界性。

\ProseParagraphBreak

球面边界有连续的一族法向方向，Fefferman 的反例正是把越来越多方向一起测试。另一方面，对一个固定方向的单个驻相片，非零曲率又能提供核衰减。分析某个新频率区域时，应同时询问局部片的衰减怎样、总共有多少片、不同片怎样组合；只有第一项回答得很好，仍可能在全局组合时失败。限制理论、多线性横截性和解耦就是为这种组合问题提供额外结构。

\ProseParagraphBreak

本章保留的必要条件与充分范围因此承担不同职责。尾部检验排除不可能的指数，可积核与解析插值提供完整可复核的充分区间，限制端点给出另一类混合指数界，深层球乘子定理解释零阶硬截断的特殊障碍。对某个尚未覆盖的正阶参数，这些结果共同界定问题所在，但没有自动给出结论。阅读研究文献中的改进时，应把其参数和模型放回这些已经明确的坐标中，才能判断它究竟改善了哪一段范围。
